\documentclass[11pt]{article}
\usepackage[a4paper,margin=0.9in]{geometry}
\usepackage{amsmath,amssymb,booktabs,graphicx,enumitem,xcolor,array}
\usepackage[hidelinks]{hyperref}
\usepackage{listings}
\usepackage{caption}
\setlength{\parindent}{0pt}\setlength{\parskip}{5pt}
\setlist{itemsep=2pt,topsep=3pt}
\renewcommand{\arraystretch}{1.2}
\newcommand{\T}{^{\top}}

\lstset{basicstyle=\ttfamily\small,breaklines=true,frame=single,framesep=4pt,
        backgroundcolor=\color{gray!8},columns=fullflexible,keepspaces=true,
        aboveskip=6pt,belowskip=6pt}

\title{\textbf{Eigenfaces vs Fisherfaces: Viva Study Guide}\\[3pt]
\large The submitted code and experiments}
\author{CS 740}
\date{}

\begin{document}
\maketitle
\vspace{-1.3em}

This guide covers only what is in the submission: the three scripts
\texttt{toy.m}, \texttt{eigenfaces\_vs\_}\allowbreak\texttt{fisherfaces.m} and
\texttt{experiments.m}, and the seven experiments E1--E7. For each piece it
gives the theory, the code that implements it, and what can be asked about it.

{\small\tableofcontents}

\vspace{6pt}\hrule

\section{The idea in one page}

Both methods turn a face image of thousands of pixels into a short list of
numbers, then recognise a new face by finding the nearest training face in that
short list. They differ in \emph{which} short list.

\textbf{PCA (Eigenfaces)} keeps the directions in which the images vary most.
It never looks at the labels. That is the whole problem: in face images the
biggest variation is often \emph{not} who the person is. Change the lighting on
one face and the pixels change far more than they do between two different
people under the same lamp. PCA faithfully keeps the lighting.

\textbf{LDA (Fisherfaces)} uses the labels. It looks for directions where the
class \emph{means} are far apart relative to the scatter \emph{within} each
class. Variation that happens inside a single person --- their own lighting,
their own expressions --- is exactly what it tries to divide out.

So the one-sentence claim of the paper is: \emph{when the nuisance variation
within a person is large, supervision is worth a great deal.} Our Yale result
(34.23\% $\to$ 67.23\%) confirms it and our ORL result (95.31\% $\to$ 90.62\%)
shows the other side: when the within-person variation is small and
unstructured, LDA's extra machinery costs more than it earns.

\section{Notation}

\begin{center}
\begin{tabular}{cll}\toprule
symbol & meaning & value in the submission\\\midrule
$d$ & pixels per image & ORL $10{,}304$, Yale $32{,}256$\\
$N$ & training images & ORL $192$, Yale $1520$\\
$c$ & classes (people) & ORL $32$, Yale $38$\\
$n_j$ & training images of person $j$ & ORL $6$, Yale $40$\\
$m$ & PCA dimension kept before LDA & swept in E1\\
$x_i$ & one image as a column & $x_i\in\mathbb R^{d}$\\
$\mu$, $\mu_j$ & overall mean, class-$j$ mean & \\
$X$ & centred training images as columns & $d\times N$\\
$U$ & eigenfaces as columns & $d\times(N-1)$\\
$W_{opt}$ & Fisherfaces as columns & $d\times(c-1)$\\
$S_T,S_B,S_W$ & total, between-class, within-class scatter & \\
$J(w)$ & Fisher criterion $\dfrac{w\T S_Bw}{w\T S_Ww}$ & \\\bottomrule
\end{tabular}
\end{center}

An image is flattened \emph{column by column} into a vector. This throws away
the fact that neighbouring pixels are adjacent, which is fine because every
image is flattened the same way, so corresponding pixels land in corresponding
coordinates.

\section{What each script does}

\begin{center}
\small
\begin{tabular}{@{}p{5.1cm}p{4.0cm}p{5.2cm}@{}}\toprule
file & runs & produces\\\midrule
\texttt{toy.m} & the 2-D demonstration & \texttt{toy.png}\\
\texttt{eigenfaces\_vs\_fisherfaces.m} & the four-method comparison & main table, two face montages\\
\texttt{experiments.m} & E1--E7 & eight figures\\\bottomrule
\end{tabular}
\end{center}

Each is a plain script with its helper functions at the bottom. MATLAB requires
a script's functions to come after all executable code, and those functions
cannot be seen by any other script --- which is why each file carries its own
copy.

\section{The toy problem (\texttt{toy.m})}

\textbf{Why it exists.} To make the difference visible in two dimensions, where
you can check the answer by eye.

Two classes share the covariance $\Sigma=\mathrm{diag}(16,\,0.25)$ and their
means differ only in $y$. So the standard deviation along $x$ is $4$ and along
$y$ is $0.5$: \emph{all the spread is along $x$ and all the class information
is along $y$}. It is a trap built for PCA.

\begin{lstlisting}
X1 = [4*randn(1, n); 0.5*randn(1, n) + 1];
X2 = [4*randn(1, n); 0.5*randn(1, n) - 1];
Sw = (X1 - m1)*(X1 - m1)' + (X2 - m2)*(X2 - m2)';
w  = Sw \ (m1 - m2);
w  = w / norm(w);
\end{lstlisting}

\textbf{The maths.} For two classes the maximiser of $J$ has a closed form.
Write $\delta=\mu_1-\mu_2$. Then $S_B=\tfrac n2\delta\delta\T$, so
$w\T S_Bw=\tfrac n2(\delta\T w)^2$, and maximising that subject to
$w\T S_Ww=1$ with a Lagrange multiplier gives
\[ 2(\delta\T w)\delta=2\lambda S_Ww\quad\Longrightarrow\quad w\propto S_W^{-1}(\mu_1-\mu_2). \]
The leading factor is a scalar, and $J$ is scale-invariant (numerator and
denominator are both quadratic in $w$), so we are free to normalise $w$.

\textbf{Result.} LDA picks $(-0.004,\,1.000)$ --- essentially the $y$-axis ---
and gets $97.5\%$. PCA picks $(-1.000,\,-0.001)$, the $x$-axis, and gets
$48.5\%$, where chance is $50\%$.

\textbf{Syntax worth knowing.} \verb|Sw \ (m1-m2)| \emph{solves}
$S_Ww=\mu_1-\mu_2$; it does not form an inverse. Backslash is a solve, not a
division, and it is both faster and more accurate than \texttt{inv(Sw)*(m1-m2)}.

\section{Eigenfaces (\texttt{eigenfaces\_vs\_fisherfaces.m})}

\subsection{The criterion}
PCA keeps the top eigenvectors of the total scatter
$S_T=\sum_i(x_i-\mu)(x_i-\mu)\T$, i.e.\ the directions of greatest variance.
Equivalently it minimises squared reconstruction error at every $k$.

\subsection{The $L=X\T X$ trick, and why it is necessary}
$S_T$ is $d\times d$. For Yale that is $32{,}256^2$ doubles, about $8.3$\,GB ---
we cannot even store it. We do not have to.

\textbf{Theorem.} Let $L=X\T X$, which is $N\times N$. If $Lw=\lambda w$ with
$\lambda\ne0$, then $Xw$ is an eigenvector of $XX\T$ with the same $\lambda$.

\textbf{Proof.} Multiply $X\T Xw=\lambda w$ on the left by $X$:
\[ X(X\T Xw)=X(\lambda w)\quad\Longrightarrow\quad (XX\T)(Xw)=\lambda(Xw), \]
regrouping by associativity. It is non-zero because
$\|Xw\|^2=w\T X\T Xw=\lambda\|w\|^2\ne0$. $\square$

So we solve a $1520\times1520$ problem instead of a $32{,}256\times32{,}256$ one.

\begin{lstlisting}
L = A' * A;
[V, D] = eig(L);
[vals, order] = sort(diag(D), 'descend');
V = V(:, order);
U = A * V;
for i = 1:N
    U(:, i) = U(:, i) / norm(U(:, i));
end
U = U(:, 1:N-1);
\end{lstlisting}

\textbf{Three things that would be wrong without a line of code.}
\begin{itemize}
\item \textbf{\texttt{eig} returns ascending eigenvalues.} Without the
  \texttt{'descend'} sort you keep the \emph{least} useful directions. It does
  not error --- the accuracy just collapses.
\item \textbf{The columns need renormalising.} The same proof gives
  $\|Xw\|=\sqrt\lambda$, not $1$. Without it $U\T U\ne I$ and the projection
  is no longer orthogonal.
\item \textbf{The last column is dropped.} After centring the columns of $X$
  sum to zero, so $\operatorname{rank}X\le N-1$ and the last eigenvalue is zero
  up to round-off. Dividing by that tiny norm would amplify pure round-off into
  something that looks like a legitimate eigenface.
\end{itemize}

\subsection{Why centre at all?}
PCA finds directions of maximum \emph{variance}, which is spread about the
centre. Without subtracting $\mu$ the first component points at the mean face
itself --- the direction in which the raw vectors are longest --- describing
what faces have in common rather than how they differ.

\subsection{The correlation baseline}
Nearest neighbour on raw pixels, after making each image zero-mean and
unit-norm. That turns Euclidean distance into correlation matching, because for
unit vectors
\[ \|a-b\|^2 = a\T a - 2a\T b + b\T b = 2 - 2\,a\T b, \]
so the nearest neighbour is the most correlated one. It also removes a
brightness offset and a contrast scaling. If a projection cannot beat this, it
is not earning its place --- and on Yale, eigenfaces do not (42.95 vs 34.23).

\section{Fisherfaces (\texttt{eigenfaces\_vs\_fisherfaces.m})}

\subsection{The criterion and the generalised eigenproblem}
\[ J(w)=\frac{w\T S_Bw}{w\T S_Ww},\qquad
   W_{fld}=\arg\max_W\frac{|W\T S_BW|}{|W\T S_WW|}. \]
Stationary points of $J$ satisfy $S_Bw=\lambda S_Ww$, and at a solution
$\lambda=J(w)$ exactly --- \emph{each generalised eigenvalue is the Fisher score
of its own direction}, which is why taking the largest ones is correct and not
merely a heuristic.

\subsection{The two rank bounds}

\textbf{$\operatorname{rank}S_B\le c-1$.} $S_B$ is a sum of $c$ rank-one terms
spanned by $\mu_j-\mu$, and those satisfy one linear relation:
\[ \sum_j n_j(\mu_j-\mu)=\sum_j n_j\mu_j-N\mu=N\mu-N\mu=0. \]
$c$ vectors with a dependence span at most $c-1$ dimensions. \textbf{Hence at
most $c-1$ Fisherfaces, however many pixels or images we have.}

\textbf{$\operatorname{rank}S_W\le N-c$.} Within class $j$ the $n_j$ centred
columns sum to zero by the definition of $\mu_j$, so they span at most $n_j-1$
dimensions; summing, $\sum_j(n_j-1)=N-c$. Read as \emph{degrees of freedom}:
each person donates only $n_j-1$ directions --- ORL 5, Yale 39.

\subsection{The small-sample-size problem}
In pixel space $S_W$ is $d\times d$ but has rank at most $N-c$. For Yale that
is $1482\ll32{,}256$: $S_W$ is massively singular, $S_W^{-1}$ does not exist,
and $J(w)$ is undefined for any $w$ in the null space. The paper's fix, which
we follow, is to project onto $m$ PCA directions first. Combining the bounds,
\[ \boxed{\,c-1\;\le\;m\;\le\;N-c\,} \]
--- the lower bound because $c-1$ directions cannot be extracted from a space
of fewer than $c-1$ dimensions. We found that by crashing the E1 sweep at
$m=20$ while asking for 31 Fisherfaces.

\subsection{Solving it}
$S_W^{-1}S_B$ is not symmetric: its transpose is $S_BS_W^{-1}$. The clean route
factors $S_W=R\T R$ (Cholesky, possible because $S_W$ is symmetric positive
definite in the reduced space) and substitutes $w=R^{-1}z$:
\[ S_BR^{-1}z=\lambda R\T RR^{-1}z=\lambda R\T z
   \ \xrightarrow{\times R^{-\top}}\
   \underbrace{R^{-\top}S_BR^{-1}}_{M}z=\lambda z, \]
and $M\T=R^{-\top}S_B\T R^{-1}=M$. So a generalised problem becomes an ordinary
\emph{symmetric} one: real eigenvalues, orthogonal eigenvectors, backward-stable
algorithm. Recover $w=R^{-1}z$; the eigenvalues are untouched, so each is still
$J(w)$.

The paper specifies none of this. It says only ``the generalized eigenvectors of
$S_B$ and $S_W$'', so the solver is our choice --- the code calls
\texttt{eig(Sb,Sw)}.

\begin{lstlisting}
[V2, D2] = eig(Sb, Sw);
[lam, order2] = sort(diag(D2), 'descend');
Wfld = V2(:, order2(1:c-1));
Wopt = Wpca * Wfld;
for i = 1:size(Wopt, 2)
    Wopt(:, i) = Wopt(:, i) / norm(Wopt(:, i));
end
\end{lstlisting}

\subsection{The normalisation step --- the most likely question}
\texttt{eig(Sb,Sw)} returns columns satisfying $W\T S_WW=I$. They are
$S_W$-orthogonal, \textbf{not} Euclidean-orthogonal, and their ordinary lengths
differ by a factor of $287$ here. $J$ does not care, because it is
scale-invariant and the \emph{subspace} is already correct. But the classifier
is 1-NN with Euclidean distance, and
\[ \|a-b\|^2=\sum_l (a_l-b_l)^2 \]
is \emph{not} scale-invariant: a coordinate whose basis vector is long
contributes more to every distance, so a few directions dominate every
comparison.

\textbf{Measured cost of omitting four lines: ORL $90.6\%\to29.5\%$.} Same
subspace, same eigenvalues, same theory --- only the basis lengths differ. This
was a real bug in our first implementation, and $29.5\%$ reads as ``LDA does not
work'', which is why it is worth being able to explain.

\subsection{Why $S_B$ needs no recentring in the code}
In the reduced space the code writes \texttt{Sb = Sb + Np*(mp*mp')} rather than
$n_j(\mu_j-\mu)(\mu_j-\mu)\T$. That is correct because $Z=W_{pca}\T X$ and $X$
was already centred, so the overall mean of $Z$ is exactly zero.

\section{The experiments}

\subsection{E1 --- sweeping $m$, with $\kappa(S_W)$}
\textbf{Question.} The paper fixes $m=N-c$, where $S_W$ has rank $m$ exactly:
full rank with \emph{zero slack}. Is that the best choice?

\textbf{Answer: no, on either database.} ORL peaks at $m=60$ with $95.3\%$
against $90.6\%$ at $m=160$; Yale at $m=1000$ with $69.8\%$ against $67.2\%$ at
$m=1482$. $\kappa(S_W)$ reaches $7.9\times10^{7}$ and $6.0\times10^{12}$
respectively at $m=N-c$.

\textbf{Mechanism.} $\kappa=\lambda_{\max}/\lambda_{\min}$ says how much a
relative error is amplified; roughly $\log_{10}\kappa$ of our 16 digits are
lost. LDA \emph{inverts} $S_W$, so it divides by the smallest eigenvalues ---
precisely the worst-estimated ones. Accuracy turns down before conditioning
becomes alarming, so the extra directions are also fitting noise, not only
losing precision.

\textbf{The deeper answer, if pushed.} The controlling quantity is not $m$ but
the sampling ratio $\gamma=m/(N-c)$. Marchenko--Pastur gives
$\kappa\approx\big((1+\sqrt\gamma)/(1-\sqrt\gamma)\big)^2$, which diverges as
$\gamma\to1$ --- and $\gamma=1$ \emph{is} $m=N-c$. Against our ORL sweep the
ratio of measurement to prediction is $9.8,\,10.1,\,8.8,\,6.5$ at
$m=31,60,100,140$: near constant, because the measured conditioning factorises
into the intrinsic conditioning of the true $S_W$ (about 9) times the sampling
inflation. So the blow-up is a predictable sampling artefact, not a quirk of
faces.

\subsection{E2 --- label permutation}
\textbf{Question.} Is $90.6\%$ real, or is the pipeline leaking?

Shuffle the training labels and change nothing else. Eigenfaces fall $94.5\to0.8$,
Fisherfaces $90.6\to0.0$, against chance $100/c=3.1\%$.

\textbf{Why below chance, not at it?} The shuffle is applied to training labels
only, so LDA faithfully separates the \emph{shuffled} grouping. A test face is
projected next to the training images that genuinely resemble it --- and those
now carry somebody else's label by construction. The method is
\emph{systematically} wrong, not randomly wrong. A method that ignored the
labels would sit at $3.1\%$; scoring below that is evidence the labels are
being used.

\subsection{E3 --- cross-validated $m$}
\textbf{Question.} E1 read ``the best $m$'' off the test column. That is tuning
on test. What is the honest number?

Six-fold leave-one-out inside the training set, with the PCA basis recomputed
per fold so the validation image cannot leak into it. CV picks $m=40$, scoring
$93.0\%$; tuning on test would pick $m=60$ for $95.3\%$.

\textbf{The 2.3-point gap is the optimism avoided.} Note the CV rates sit
\emph{above} the test rates, which is consistent rather than contradictory: the
validation images come from the same six-image pool as the training images, an
easier problem. CV is used to \emph{rank} values of $m$, not to estimate the
final accuracy.

\subsection{E4 --- noise robustness}
\textbf{A prediction, not a fishing trip.} Fisherfaces are higher-frequency
than eigenfaces because whitening scales each eigen-direction of $S_W$ by
$1/\sqrt\lambda$, amplifying the low-variance fine-detail directions most.
White noise has energy at every frequency, so Fisherfaces should degrade faster.

\textbf{Confirmed, with a caveat.} From $\sigma=0$ to $640$ grey levels
Fisherfaces lose $60.2$ points against eigenfaces' $21.9$. But nothing moves
below $\sigma\approx160$: our first attempt stopped at $\sigma=80$ and found
all methods moving $0.8$ points, which on 128 test images is exactly
\emph{one image}. Projection is a strong denoiser --- a unit-norm direction
turns per-pixel noise of standard deviation $\sigma$ into coefficient noise of
the same $\sigma$, while the coefficients themselves are in the thousands.

\subsection{E5 --- error against lighting angle}
Yale filenames encode the lamp direction, e.g.\
\texttt{yaleB01\_P00A+020E-40} is azimuth $+20^\circ$, elevation $-40^\circ$.
Treating them as rotations about perpendicular axes, the direction cosine along
the camera axis is the product of the cosines, so
$\theta=\arccos(\cos\mathrm{az}\cdot\cos\mathrm{el})$.

\begin{center}
\begin{tabular}{lccccc}\toprule
Error (\%) & $0$--$12^\circ$ & $13$--$25^\circ$ & $26$--$50^\circ$ & $51$--$77^\circ$ & $>77^\circ$\\
Images & 0 & 104 & 152 & 300 & 338\\\midrule
Eigenfaces & --- & 1.0 & 30.9 & 88.3 & 93.8\\
Fisherfaces & --- & 0.0 & 2.0 & 22.0 & 66.3\\\bottomrule
\end{tabular}
\end{center}

Error climbs and the gap \emph{widens}: $15\times$ better at $26$--$50^\circ$.

\textbf{The empty first bin is itself a finding.} \texttt{dir} sorts names as
text and \texttt{'+'} precedes \texttt{'-'}, so taking the first 40 files per
subject puts every near-frontal lamp into training. Every test image is at
least $12^\circ$ off axis and 638 of 894 are beyond $50^\circ$. The test set is
the hard tail of the lighting distribution, not a sample of it --- which is why
the headline Yale numbers look so poor.

\subsection{E6 --- sweeping the number of classes}
For $c=2,4,8,16,32$ the number of Fisherfaces produced is $1,3,7,15,31$ ---
exactly $c-1$ every time, so the rank bound is a live constraint and not merely
a proof on a page. Fisherfaces are at $100\%$ for $c\le8$ on the \emph{same}
six-images-per-person data on which they lose at $c=32$, which is why ``too few
images per class'' cannot on its own be the explanation for ORL.

\subsection{E7 --- glasses or no glasses}
A different question: not \emph{who} is this but \emph{is this person wearing
glasses}. With $c=2$, $\operatorname{rank}S_B\le1$ and there is exactly
\textbf{one} Fisherface, $w\propto S_W^{-1}(\mu_1-\mu_2)$. PCA's subspace
depends only on $S_T$ and does not change with the labels at all --- which is
the cleanest illustration of supervised versus unsupervised in the whole
project.

ORL ships no glasses labels, so all 400 images were labelled by inspection: 117
wear glasses, 8 people appear both ways. On the raw set, always answering
``no'' scores $70.8\%$ and beats any real method, so the set was balanced to
234 images and evaluated leave-one-\emph{subject}-out, so the test face is never
seen in training.

\textbf{Result: PCA $44.9\%$ error, one Fisherface $41.0\%$, baseline $50\%$.}
The paper got $52.6\%$ and $5.3\%$. \emph{We reproduce that PCA is near chance
but not that Fisherfaces solve the task.} The likely cause is the data: the
paper's 36 images were \emph{paired} --- the same person with and without
glasses in one session --- so glasses were the only thing that changed. In ORL
glasses co-vary with pose, expression and session, so one direction has no
clean signal to lock onto.

\section{MATLAB constructs used}

\begin{center}
\begin{tabular}{ll}\toprule
construct & what to say about it\\\midrule
\texttt{im(:)} & flattens an image to a column, \emph{column by column}\\
\texttt{reshape(v,h,w)} & the exact inverse; used to display a face\\
\verb|A \ b| & \emph{solves} $Ax=b$; never forms an inverse\\
\texttt{eig(L)} & returns eigenvalues \textbf{ascending}; we always sort\\
\texttt{eig(Sb,Sw)} & the generalised problem; columns are $S_W$-orthonormal\\
\texttt{sort(\dots,'descend')} & second output is the permutation; reorder vectors with it\\
\texttt{double(imread(\dots))} & \texttt{uint8} arithmetic saturates --- convert first\\
\texttt{X - meanFace} & implicit expansion subtracts from all columns\\
\texttt{Z(:, labels==p)} & logical indexing picks one class's columns\\
\texttt{cond(Sw)} & $\lambda_{\max}/\lambda_{\min}$; E1 plots it\\
\texttt{rank(Sw)} & tolerance-based, so it can disagree with \texttt{chol}\\\bottomrule
\end{tabular}
\end{center}

\section{Viva questions}

\subsection{Core}
\begin{description}[style=nextline,leftmargin=0pt]
\item[Why does PCA fail under lighting?] It keeps directions of largest
  variance and never sees the labels. Changing the lamp moves the pixels more
  than changing the person does, so the leading components encode lighting.
\item[Then why not just drop the first three components?] You can, and it
  helps: Yale goes $34.23\to58.72$. But it assumes the nuisance is exactly
  3-dimensional (the Lambertian argument), whereas LDA learns the nuisance from
  the labels. Fisherfaces still beat drop-3 on Yale. And on ORL the same trick
  \emph{hurts} ($95.31\to87.50$), because there the leading components carry
  identity.
\item[Why do Fisherfaces lose on ORL?] Eigenfaces never estimates $S_W$; it
  keeps the \emph{largest} eigenvalues of a well-estimated $S_T$. Fisherfaces
  must \emph{invert} $S_W$ and so leans on its \emph{smallest} eigenvalues,
  which are the worst estimated. ORL gives 5 degrees of freedom per person,
  Yale 39.
\item[So it is just a sample-size problem?] Not only. E6 shows Fisherfaces at
  $100\%$ for $c\le8$ on the same six-images-per-person data. What matters is
  how much \emph{structured} within-class variation there is to remove ---
  illumination on Yale is large, shared and nearly low-dimensional; pose and
  expression on ORL are small and idiosyncratic.
\end{description}

\subsection{Linear algebra}
\begin{description}[style=nextline,leftmargin=0pt]
\item[Why is there a limit of $c-1$ Fisherfaces?] $\operatorname{rank}S_B\le c-1$,
  because $\sum_jn_j(\mu_j-\mu)=0$ makes the class-mean vectors dependent.
  E6 confirms the count is exactly $c-1$ at every $c$ tested.
\item[Why is $S_W$ singular in pixel space?] $\operatorname{rank}S_W\le N-c$,
  since each class's deviations sum to zero and contribute $n_j-1$ directions.
  For Yale $1482\ll32{,}256$.
\item[Why not just use \texttt{inv(Sw)*Sb}?] It is not symmetric --- its
  transpose is $S_BS_W^{-1}$ --- so you lose the guarantees of a symmetric
  solver: real eigenvalues and orthogonal eigenvectors.
\item[Why does the Cholesky substitution work?] $S_W=R\T R$, and $w=R^{-1}z$
  turns $S_Bw=\lambda S_Ww$ into $(R^{-\top}S_BR^{-1})z=\lambda z$, which is
  symmetric because $S_B$ is. Eigenvalues are unchanged, so each is still $J(w)$.
\item[Are Fisherfaces orthogonal?] \textbf{No.} They are $S_W$-orthogonal:
  $W\T S_WW=I$, not $W\T W=I$. Measured, the off-diagonal energy is
  $3.8\times10^{-13}$ in the $S_W$ inner product and $0.499$ in the Euclidean
  one --- half the energy off-diagonal.
\item[Why is $\|Xw\|=\sqrt\lambda$?] $\|Xw\|^2=w\T X\T Xw=w\T(\lambda w)=\lambda$
  for unit $w$. Hence the renormalisation loop.
\end{description}

\subsection{Implementation}
\begin{description}[style=nextline,leftmargin=0pt]
\item[What breaks if you skip the unit-norm loop on \texttt{Wopt}?] ORL falls
  from $90.6\%$ to $29.5\%$. \texttt{eig(Sb,Sw)} gives $S_W$-orthonormal
  columns whose Euclidean lengths differ by $287\times$, and 1-NN with
  Euclidean distance is not scale-invariant.
\item[What breaks if you forget \texttt{'descend'}?] Nothing errors --- you
  silently keep the least useful directions and accuracy collapses.
\item[Why drop the last eigenface?] After centring $\operatorname{rank}X\le N-1$,
  so the last eigenvalue is zero up to round-off; dividing by that norm turns
  round-off into a unit vector that looks legitimate.
\item[Why \texttt{double()} before any arithmetic?] \texttt{uint8} saturates:
  $200+100$ becomes $255$ and $5-10$ becomes $0$, so every mean and difference
  would be silently wrong.
\item[Why is $m=N-c$ and not larger?] Larger makes $S_W$ singular, since
  $\operatorname{rank}S_W\le N-c$. And not smaller than $c-1$, since you cannot
  pull $c-1$ directions out of fewer dimensions.
\item[How do you know your accuracy is not a bug?] E2: destroying the training
  labels takes both methods to chance or below.
\end{description}

\subsection{Protocol and honesty}
\begin{description}[style=nextline,leftmargin=0pt]
\item[Is $95.3\%$ your real ORL accuracy?] No. That is the best $m$ chosen by
  looking at the test set. Cross-validated, $m=40$ gives $93.0\%$, and the
  2.3-point gap is the optimism. $93.0\%$ is the honest figure.
\item[Why are the Yale numbers so low?] The split. \texttt{dir} sorts as text
  and \texttt{'+'} precedes \texttt{'-'}, so training gets every frontal lamp
  and testing gets only hard angles --- every test image is at least $12^\circ$
  off axis and 638 of 894 beyond $50^\circ$.
\item[What did you not implement?] The Linear Subspace method of the paper's
  Section 2.3, which builds a 3-D illumination subspace per person.
\item[Did you use the paper's databases?] No, and the report says so. The
  paper's lighting experiment used the Harvard database, which has no surviving
  public download, and its Yale set was a smaller 1997 database. We used
  Cropped Extended Yale B, by the same authors in 2001, so E5 is modelled on
  the paper's protocol rather than a reproduction of it.
\item[Anything that did not work?] Yes, and it is reported: E7. We reproduce
  the paper's finding that PCA is near chance on glasses, but not that
  Fisherfaces solve it ($41.0\%$ error against their $5.3\%$). Their images were
  paired, same person with and without glasses in one session; ORL's are not.
\end{description}

\section{One-glance summary}

\begin{center}
\begin{tabular}{ll}\toprule
Fact & Number\\\midrule
Yale: eigenfaces $\to$ Fisherfaces & $34.23\% \to 67.23\%$\\
ORL: eigenfaces $\to$ Fisherfaces & $95.31\% \to 90.62\%$\\
Drop top 3 on Yale (helps) & $34.23\% \to 58.72\%$\\
Drop top 3 on ORL (hurts) & $95.31\% \to 87.50\%$\\
Fisherfaces without unit-norming & $90.6\% \to 29.5\%$\\
Honest ORL accuracy (cross-validated) & $93.0\%$, not $95.3\%$\\
Best $m$ vs the paper's $N-c$ (ORL) & $60$ vs $160$; $95.3$ vs $90.6$\\
$\kappa(S_W)$ at $m=N-c$ & $7.9\times10^{7}$ (ORL), $6.0\times10^{12}$ (Yale)\\
Number of Fisherfaces & exactly $c-1$, every $c$ tested\\
Labels shuffled & $0.8\%$ and $0.0\%$ vs chance $3.1\%$\\
Lighting $26$--$50^\circ$ & $30.9\%$ vs $2.0\%$ error\\
Glasses, one Fisherface & $41.0\%$ error vs the paper's $5.3\%$\\\bottomrule
\end{tabular}
\end{center}

\end{document}
